\chapter{Hilbert--Schmidt Geometry and the Frobenius Norm}
\label{app:hilbert-schmidt-and-frobenius}

\textbf{Status: expository appendix.} This appendix records the operator-space geometry
underlying Frobenius-norm comparisons of reduced operators. It does not define
a scientific acceptance metric.

\section{Hilbert--Schmidt inner product}

Let $V$ be a finite-dimensional complex Hilbert space, and let
$\mathcal L(V)$ denote the vector space of linear operators from $V$ to itself.
If $\dim V=d$, then $\dim\mathcal L(V)=d^2$. The function
\begin{equation}
    \langle\cdot,\cdot\rangle_{\mathrm{HS}}:
    \mathcal L(V)\times\mathcal L(V)
    \longrightarrow
    \mathbb C
\end{equation}
defined by
\begin{equation}
    \langle A,B\rangle_{\mathrm{HS}}
    =
    \operatorname{Tr}\!\left(A^\dagger B\right)
\end{equation}
is the \emph{Hilbert--Schmidt inner product}, also called the trace inner
product~\cite{nielsenchuang2010}. With this convention it is conjugate-linear
in its first argument and linear in its second. Its induced norm is
\begin{equation}
    \lVert A\rVert_{\mathrm{HS}}
    =
    \sqrt{
        \operatorname{Tr}\!\left(A^\dagger A\right)
    }.
\end{equation}

Let $\mathcal H^{(P)}$ be a finite-dimensional retained Hilbert space. The
operator space
\begin{equation}
    \mathcal L\!\left(\mathcal H^{(P)}\right)
\end{equation}
is itself a finite-dimensional Hilbert space under this inner product. If
\begin{equation}
    \hat A,\hat B:
    \mathcal H^{(P)}
    \longrightarrow
    \mathcal H^{(P)},
\end{equation}
then their sum and difference are well-defined elements of the same operator
space:
\begin{equation}
    \hat A\pm\hat B:
    \mathcal H^{(P)}
    \longrightarrow
    \mathcal H^{(P)}.
\end{equation}

Given an orthonormal basis of $\mathcal H^{(P)}$, let $\mathbf A$ and
$\mathbf B$ be the matrices representing $\hat A$ and $\hat B$. The
Hilbert--Schmidt inner product becomes the Frobenius inner product,
\begin{equation}
    \langle\mathbf A,\mathbf B\rangle_{\mathrm F}
    =
    \operatorname{Tr}\!\left(\mathbf A^\dagger\mathbf B\right),
\end{equation}
and the corresponding norm satisfies
\begin{equation}
    \lVert\mathbf A\rVert_{\mathrm F}^2
    =
    \operatorname{Tr}\!\left(\mathbf A^\dagger\mathbf A\right)
    =
    \sum_{i,j}\lvert A_{ij}\rvert^2.
\end{equation}
Thus the Frobenius norm is the matrix-coordinate form of the
Hilbert--Schmidt norm when the coordinates are defined by an orthonormal basis.

Operators may be added and subtracted because they belong to a common operator
vector space. The legitimacy of an operator subtraction does not follow merely
from two matrices having the same dimensions. It follows from the represented
operators having compatible domains and codomains and from the matrices using
an identified common state space and basis convention.

\section{Common projection and reduced components}

Let $\hat P$ be an orthogonal projector and define
\begin{equation}
    \mathcal H^{(P)}
    =
    \operatorname{Ran}\hat P.
\end{equation}
For possibly unbounded parent operators $\hat H$ and $\hat T$, assume that the
retained range lies in the domains needed to define the compressions below.
Set
\begin{align}
    \hat H_{\mathrm{eff}}
    &=
    \left.
        \hat P\hat H\hat P
    \right|_{\mathcal H^{(P)}},
    \\
    \hat T_{\mathrm{eff}}
    &=
    \left.
        \hat P\hat T\hat P
    \right|_{\mathcal H^{(P)}}.
\end{align}
Both reduced operators belong to
$\mathcal L\!\left(\mathcal H^{(P)}\right)$. Their difference is therefore a
well-defined reduced operator,
\begin{equation}
    \hat V_{\mathrm{eff}}
    =
    \hat H_{\mathrm{eff}}
    -
    \hat T_{\mathrm{eff}},
\end{equation}
and, by construction,
\begin{equation}
    \hat H_{\mathrm{eff}}
    =
    \hat T_{\mathrm{eff}}
    +
    \hat V_{\mathrm{eff}}.
\end{equation}

If the parent operators satisfy
$\hat H=\hat T+\hat V$ on a compatible domain, linearity gives
\begin{equation}
    \hat V_{\mathrm{eff}}
    =
    \left.
        \hat P\hat V\hat P
    \right|_{\mathcal H^{(P)}}.
\end{equation}
Even when $\hat V$ is a local multiplication operator in the parent position
representation, its compression need not act as multiplication by a scalar
function within the retained representation. Projection can produce an
operator that is nonlocal in the retained coordinates. If the parent
Hamiltonian contains nonlocal pseudopotential, spin--orbit, or other operator
components, the chosen kinetic--remainder decomposition must identify those
components explicitly.

\section{Aligned pristine--doped operators}

The same common-space requirement applies to independently constructed
pristine and doped retained operators. Suppose first that the two retained
spaces have equal finite dimension and that a physically admissible unitary
identification has been constructed,
\begin{equation}
    \hat U_d:
    \mathcal H_b^{(P)}
    \longrightarrow
    \mathcal H_d^{(P)}.
\end{equation}
Let $c_b$ and $c_d$ denote scalar energy references selected by the declared
alignment procedure, and define
\begin{equation}
    \overline{\hat H}_s^{(P)}
    =
    \hat H_s^{(P)}-c_s\hat I_s,
    \qquad
    s\in\{b,d\}.
\end{equation}
The aligned impurity operator on the bulk retained space is then
\begin{equation}
    \Delta\hat H_d^{(P)}
    =
    \hat U_d^\dagger
    \overline{\hat H}_d^{(P)}
    \hat U_d
    -
    \overline{\hat H}_b^{(P)}
    \in
    \mathcal L\!\left(\mathcal H_b^{(P)}\right).
\end{equation}
Equal matrix dimensions alone would not establish this membership. The
identification must also preserve the physical structures required by the
comparison, including the applicable gauge, geometry, orbital ordering, spin
structure, and energy reference.

The Hilbert--Schmidt discrepancy of the aligned operator is
\begin{equation}
    \left\lVert
        \Delta\hat H_d^{(P)}
    \right\rVert_{\mathrm{HS}}^2
    =
    \operatorname{Tr}\!\left[
        \left(\Delta\hat H_d^{(P)}\right)^\dagger
        \Delta\hat H_d^{(P)}
    \right].
\end{equation}
In a common orthonormal numerical basis, this is exactly
\begin{equation}
    \left\lVert
        \Delta\mathbf H_d^{(P)}
    \right\rVert_{\mathrm F}^2
    =
    \sum_{\alpha,\beta}
    \left\lvert
        \Delta H_{d,\alpha\beta}^{(P)}
    \right\rvert^2.
\end{equation}
Frobenius-norm minimization is therefore least-squares approximation in the
finite-dimensional Hilbert space of retained operators.

\section{Domains and the particle in a box}

The differential expression
\begin{equation}
    -\frac{\hbar^2}{2m}\frac{d^2}{dx^2}
\end{equation}
does not by itself specify a quantum-mechanical operator. An operator is
determined by both its action and its domain. The free kinetic operator on the
real line and the kinetic operator on a finite interval with Dirichlet boundary
conditions are therefore different operators even though they share the same
differential expression.

For the particle in a box, let
\begin{equation}
    \hat H_{\mathrm{box}}
    =
    -\frac{\hbar^2}{2m}\frac{d^2}{dx^2},
    \qquad
    \mathcal D(\hat H_{\mathrm{box}})
    =
    H^2(0,L)\cap H_0^1(0,L).
\end{equation}
If $\hat T$ is this same Dirichlet kinetic operator, then
$\hat H_{\mathrm{box}}=\hat T$ and consistent projection gives
\begin{equation}
    \hat H_{\mathrm{eff}}-\hat T_{\mathrm{eff}}
    =
    \hat 0.
\end{equation}

A nonzero difference can arise if the chosen reference kinetic operator does
not carry the same confinement or boundary-domain structure. Such a comparison
first requires an explicit common-space identification. The resulting residual
can then record the missing confinement or boundary contribution relative to
the chosen reference kinetic operator; it must not be identified automatically
with a local physical potential. This distinction separates a mathematically
valid represented subtraction from its physical interpretation.

\section{Infinite-dimensional qualification}

For an infinite-dimensional Hilbert space $\mathcal H$, the space of all
bounded linear operators is not a Hilbert space under the Hilbert--Schmidt inner
product. One instead uses the Hilbert--Schmidt class
\begin{equation}
    \mathcal L_2(\mathcal H)
    =
    \left\{
        \hat A:
        \operatorname{Tr}\!\left(\hat A^\dagger\hat A\right)
        <
        \infty
    \right\}.
\end{equation}
For $\hat A\in\mathcal L_2(\mathcal H)$,
\begin{equation}
    \lVert\hat A\rVert_{\mathrm{HS}}^2
    =
    \operatorname{Tr}\!\left(\hat A^\dagger\hat A\right).
\end{equation}
Finite-rank operators are Hilbert--Schmidt. Consequently, a well-defined
compression to a finite-dimensional retained space has finite
Hilbert--Schmidt norm.

The Frobenius norm is invariant under a common unitary change of orthonormal
coordinates:
\begin{equation}
    \left\lVert
        \mathbf U^\dagger\mathbf A\mathbf U
        -
        \mathbf U^\dagger\mathbf B\mathbf U
    \right\rVert_{\mathrm F}
    =
    \lVert\mathbf A-\mathbf B\rVert_{\mathrm F}.
\end{equation}
It is not meaningful to invoke this invariance when only one operand has been
transformed or when the two matrices have not first been placed in a common
represented space.

The Hilbert--Schmidt geometry supplies a rigorous finite-dimensional setting
for operator comparison and Frobenius-norm minimization. It does not by itself
make the Frobenius norm the scientifically appropriate loss. The common state
space, retained subspace, normalization, alignment, weighting, domains, energy
reference, and relation of the norm to the observables of interest must still
be specified.
